La lumière d'une lampe à décharge contient deux raies, de longueurs d'onde \qty{\v{lam1}}{nm} et \qty{\v{lam2}}{nm}. Elle éclaire sous incidence normale un réseau de \num{\v{Nmm}} traits par millimètre. Un écran plat, parallèle au réseau, est placé à \qty{\v{L}}{m}.

\begin{enumerate}[label=\alph*)]
    \item Déterminez l'angle de chacune des deux raies dans le spectre d'ordre 1.
    \item Déterminez la position de chacune des raies sur l'écran, mesurée à partir du maximum central.
    \item Quelle distance sépare les deux raies sur l'écran ?
\end{enumerate}
